\section{Projection, matching, and the cycle count}\label{sec:counting}

We first prove a finite-dimensional counting principle with its full
closure hypothesis.  We then construct the analytic systems to which it
applies.  Throughout this section, a parameter becomes an unknown when
an absolute zero test is made.  The distinction between these two roles
is essential.

Uniform component bounds for relatively compact subanalytic families
are classical \citep[Corollary~1]{Gabrielov1968}; see also
\citet[Theorem~3.14]{BierstoneMilman1988}.
In the Pfaffian setting, Khovanskii similarly reduces component bounds
to critical-point equations of proper functions
\citep[\S4, Theorems~3--4]{Khovanskii1984Finiteness}.
Here we prove the required implication for open analytic systems with
an explicit differential and algebraic closure hypothesis, using proper
functions and Sard's theorem. The passage constructions must verify
that hypothesis before the component bound can be applied.

\subsection{Open domains and the projection principle}

\begin{definition}[Admissible open data]\label{def:counting-domain}
An open analytic system consists of a specified open set
$U\subset\RR^a\times\RR^n$, real analytic equations
$F:U\longrightarrow\RR^m$, and a finite list of real analytic margins
$g_1,\ldots,g_s>0$ specifying its open restrictions.  We require the
following completeness of the list of margins: if a sequence in $U$ has
bounded coordinates and all its margins are bounded below by positive
constants, every finite limit of that sequence belongs to $U$ and is a
point of analyticity of all the data.  Coordinate-chart boundaries are
included among these restrictions.  A nonzero condition $h\ne0$ can be
written as the margin $h^2>0$.

We call this condition \emph{positive-margin completeness}.  It concerns
bounded limits for which every listed margin stays bounded away from
zero; it does not assert analytic extension when a margin tends to zero.

The \emph{finite differential and algebraic closure} of these data allows
fixed real constants, additional real coordinates, finite sums and
products, finitely many independent-coordinate derivatives, and inverses
of functions declared nonzero on the domain.  One may clear these
denominators while retaining their nonzero restrictions.  A system used
below retains the original equations, argument graph ties, and domain
restrictions, and can append equations and margins from this closure.
The operation can be iterated finitely many times.  Fixed real values
may be assigned to designated parameters at any stage.  All such
specialized equations and restrictions are retained.
Equivalently, a specialization retains those parameters as coordinates
and appends the polynomial equations fixing their values.

We say that the data have the \emph{absolute zero property for this
closure} if every fixed finite system so obtained has only finitely
many isolated real solutions in its specified open domain, when all
remaining coordinates are treated as unknowns.  This assertion concerns
the entire system, including every added coordinate.  It is not an
assertion merely about isolated points of the original zero set.
\end{definition}

For a system with full row rank in $x$, one can use separate charts for
the finitely many nonzero $m\times m$ minors of $F_x$.  Their squares
are added to the margin list.  Every further determinant restriction
in the following proof is treated in the same way.

\begin{theorem}[Projection count]\label{thm:projection-count}
Let $F(p,x)=0$ be admissible open data as in
Definition~\ref{def:counting-domain}, with
$p\in\RR^a$, $x\in\RR^n$, and $m\le n$.  Suppose that $F_x$ has
row rank $m$ throughout the zero set under consideration and that the
data have the absolute zero property for their finite differential and
algebraic closure.  Then there is a finite constant $N$, depending on
this fixed system and its domain, such that every fiber
\[
 M_p=\{x:(p,x)\in U,\ F(p,x)=0\}
\]
has at most $N$ connected components.  If $m=n$, every fiber has at most
$N$ points.
\end{theorem}

\begin{proof}
We first prove the square case by induction on the number of parameters.
With the other parameters fixed, allowing one parameter to vary turns
its solutions into curves.  A fixed generic tilt associates their
components with a square critical-point system having one fewer
parameter.  We then treat positive-dimensional fibers by choosing a
tilt for a countable family of witness fibers and applying the square
case to their multiplier systems.

We may work on one minor chart and finally sum over the finitely many
charts.  Include its determinant margin in the list $g$.  On a fiber,
and also on a slice in which only some parameters have been fixed, use
\begin{equation}\label{eq:counting-barrier}
 \rho(p,x)=1+|p|^2+|x|^2+\sum_{j=1}^s g_j(p,x)^{-2}.
\end{equation}
Its sublevels on the closed equation locus are compact.  Indeed the
coordinates are bounded there, each $g_j$ is bounded below by a
positive constant, and Definition~\ref{def:counting-domain} puts every
limit back in the same domain; continuity then preserves $F=0$.
For any fixed linear function $\ell\cdot y$ of the varying coordinates,
\[
 |y|^2+\ell\cdot y\ge \tfrac12|y|^2-\tfrac12|\ell|^2.
\]
Consequently $\rho+\ell\cdot y$ is proper and bounded below as well.
Every nonempty connected component of a smooth fiber is closed in that
fiber and therefore contains a minimum of this function.  The minimum
is an interior constrained critical point.

We prove the assertion for square systems first, simultaneously for
all systems in the stated closure, by induction on the number $a$ of
parameters.  If $a=0$, invertibility of $F_x$ makes every solution
isolated, and the absolute zero property is exactly the required
finiteness.  Suppose $a>0$ and write $p=(b,t)$ with
$b\in\RR^{a-1}$ and $t\in\RR$.  At fixed $b$, the equation locus is
a one-dimensional analytic manifold in $y=(t,x)$, since $F_x$ is
invertible.  Its tangent vectors satisfy
\[
 F_t\,\dd t+F_x\,\dd x=0,
 \qquad \dd x=-F_x^{-1}F_t\,\dd t.
\]
Thus $\dd t$ never vanishes on a nonzero tangent vector.  The function
$t$ is injective on every connected component: a connected
one-dimensional manifold without boundary is a circle or an interval;
a circle would force an extremum of $t$, and on an interval its
nonzero derivative has a constant sign.  It follows that a fiber over
$(b,t)$ contains at most one point from each of these curve components.

We next choose one tilt, independent of $b$.  Introduce
$\lambda\in\RR^n$ and $\ell\in\RR^{n+1}$, and consider
\begin{equation}\label{eq:counting-critical}
 F(b,y)=0,
 \qquad
 \nabla_y\rho(b,y)+\ell-F_y(b,y)^{\mathsf T}\lambda=0.
\end{equation}
The universal critical locus is a smooth manifold: first impose
$F=0$, whose derivative in $x$ is surjective, and then solve the second
block for $\ell$.  Its dimension equals that of $(b,\ell)$.  The
kernel of the derivative of projection to $(b,\ell)$ is exactly the
kernel of the square Jacobian of the left sides of
\eqref{eq:counting-critical} in $(y,\lambda)$.
Sard's theorem, applied on countably many charts if necessary, says
that the critical values of this projection form a null set
\cite{Sard1942}.  Fubini's theorem therefore gives a fixed $\ell$ for
which, for almost every $b$, every solution of
\eqref{eq:counting-critical} has an invertible Jacobian in
$(y,\lambda)$.

For clarity, this invertibility is precisely nondegeneracy of a
critical point of the restricted tilted function.  If $A=F_y$ and
$H=D_y^2(\rho+\ell\cdot y-\sum_i\lambda_iF_i)$ is the ambient
Lagrangian Hessian, its restriction to $\ker A$ is the constrained
Hessian.  The relevant bordered matrix is
\[
 \begin{pmatrix} A&0\\ H&-A^{\mathsf T}\end{pmatrix}.
\]
Its kernel consists of $(v,\eta)$ with $Av=0$ and
$Hv=A^{\mathsf T}\eta$.  It is trivial exactly when the quadratic
form $H$ on $\ker A$ is nondegenerate: testing the second equation
against vectors in $\ker A$ proves one direction; the orthogonal
complement identity $(\ker A)^\perp=\operatorname{im} A^{\mathsf T}$
proves the other.

Now fix this $\ell$.  On the open locus where its displayed Jacobian
is invertible, \eqref{eq:counting-critical} is a square system with
$2n+1$ unknowns $(y,\lambda)$ and only $a-1$ parameters $b$.
It belongs to the closure in Definition~\ref{def:counting-domain}.
In fact the derivatives of $g_j^{-2}$ are products of derivatives of
$g_j$ and powers of its nonzero inverse; determinant conditions are
polynomials in such derivatives.  All original equations and
restrictions are retained.  In subsequent proper functions we include
the squares of all new coordinates, including $\lambda$, and inverse
barriers for all new strict conditions.  The induction hypothesis
therefore gives a uniform bound $N_0$ on the regular solutions of this
new system for every $b$.

For almost every $b$, every curve component has a nondegenerate
minimum and hence gives a different solution of the new system.
There are at most $N_0$ components and therefore at most $N_0$ points
in each original fiber over such a $b$.  This bound holds at an
exceptional $(b_0,t_0)$ too.  Otherwise select $N_0+1$ distinct points
of that fiber.  By the implicit function theorem they extend to
disjoint local solution graphs over one common neighborhood of
$(b_0,t_0)$, preserving every strict inequality.  Choose $b$ in the
full-measure set close to $b_0$, with $t=t_0$.  All selected points
persist, a contradiction.  Only finitely many points were selected,
so no common neighborhood for an infinite fiber was assumed.  This
completes the induction for square systems.

Finally allow $m<n$.  If no uniform component bound existed, choose
a countable collection of parameters $p_k$ with more than $k$
components in $M_{p_k}$.  For each fixed $p_k$ the same universal
critical construction, now in $(x,\lambda,\ell)$ with
$\lambda\in\RR^m$ and $\ell\in\RR^n$, shows that almost every
tilt makes $\rho(p_k,x)+\ell\cdot x$ Morse on the whole fiber.
Choose one tilt in the intersection of these countably many
full-measure sets.  The multiplier equations
\[
 F(p,x)=0,
 \qquad \nabla_x\rho(p,x)+\ell-F_x(p,x)^{\mathsf T}\lambda=0
\]
are square in $(x,\lambda)$.  Restricting to their invertible
Jacobian locus gives a system covered by the square assertion just
proved.  Each component in every witness fiber contributes a regular
minimum, contradicting its uniform point bound.  The same argument
includes $m=0$, with no multiplier coordinates.  This proves the
theorem.
\end{proof}

\begin{remark}\label{rem:counting-closure}
The closure hypothesis is hereditary, and is used each time the
parameter dimension decreases.  For example, the equation
$\sin x=0$ on $0<x<p$ has no isolated zeros in the full $(p,x)$ space,
while its finite fibers have unbounded size.  Appending $p=x+1$
produces infinitely many isolated total zeros.  Thus finiteness for
the original total zero set alone cannot replace the hypothesis of
Theorem~\ref{thm:projection-count}.  Conversely, its proof uses only
finitely many derivative orders at each of finitely many induction
steps; it imposes no estimate uniform over all derivative orders.
\end{remark}

\subsection{Matching equations and their geometric meaning}

\begin{theorem}[Matching systems]\label{thm:matching-system}
For each polynomial degree bound $d$, there are a finite number $T(d)$
and finitely many fixed open analytic systems
\[
 F_w(p_w,x_w)=0,\qquad g_w(p_w,x_w)>0,
 \qquad w\in\mathcal W_d,
\]
with the following properties.
\begin{enumerate}
\item Their domains satisfy Definition~\ref{def:counting-domain}.
Every finite differential and algebraic system generated from these
data, with their original ties and domain restrictions retained, has
the absolute zero property.
\item For each field $V=(P,Q)$ of degree at most $d$, all but at most
$T(d)$ of its periodic orbits contained in the fixed square of
Section~\ref{sec:subdivision} have a representation in one of these
systems.  The parameter value $p_w(V)$ is the same for every orbit
using the same word and preparation choices; endpoint-dependent
quantities are unknowns with graph ties.
\item At a representation of a hyperbolic periodic orbit, $F_{w,x_w}$
has full row rank.  In each fixed nonzero-minor chart, different
represented hyperbolic orbits lie on different connected components
of the fiber at $p_w(V)$.
\end{enumerate}
The minor charts and their additional margins also satisfy the first
assertion.
\end{theorem}

The proof occupies this subsection and the next.  We first construct the
systems and show why their regular fiber components distinguish hyperbolic
cycles.  Moving the endpoint cuts along an orbit gives continuous families of
representations.  After
establishing the geometric interpretation, we prove the absolute zero
property needed to bound its components.

\paragraph{Finite words and connectors.}
Theorem~\ref{thm:finite-templates} gives a finite alphabet of prepared
passages and a degree-dependent bound on the number of boundary
events of a periodic Jordan curve.  Orbits containing one of its
nonsingular tangential boundary arcs account for at most $T(d)$
exceptions.  At every other boundary event insert a small connector,
with its two ends in the interiors of the neighboring pieces.
There are only finitely many events on the selected orbit, so these
connectors can be made disjoint.  A periodic orbit cannot remain
entirely inside a piece with a nonsingular monotone graph clock.
In particular, a nonconstant closed curve has an extremum of its
horizontal coordinate.  Thus the retained cycles are encoded by
closed words of bounded length.

Here is a fixed analytic connector model.  Choose one nonzero physical
velocity component as graph clock $X$, let $Y$ be the other coordinate,
and order the ends so that $\Delta=X_+-X_->0$.  Physical time may run
in either direction.  Introduce $r>0$ by $r^2=\Delta$ and set
\[
 X=X_-+\Delta s,\qquad Y=Y_-+rz,\qquad 0\le s\le1.
\]
Writing $V_X,V_Y$ for the two components, the scalar equation is
\begin{equation}\label{eq:counting-connector}
 z_s=\frac{\widehat N(s,z)}{\widehat D(s,z)},\qquad
 \widehat D=\frac{V_X(X,Y)}{V_X(X_-,Y_-)},\qquad
 \widehat N=\frac{\Delta}{r}
              \frac{V_Y(X,Y)}{V_X(X_-,Y_-)}.
\end{equation}
These are polynomials of degree bounded in terms of $d$ in $(s,z)$.
Their coefficients are tied to the original field and the endpoints
by polynomial equations, with the divisor $V_X(X_-,Y_-)$ retained
as nonzero.  As the connector is shortened at a fixed nonsingular
event, the coefficient vectors of $(\widehat D,\widehat N)$ tend to
$(1,0)$.  Indeed every nonconstant term of $\widehat D$ contains a
factor $\Delta$ or $r$, and every term of $\widehat N$ contains the
factor $\Delta/r=r$.  Also $z_-=0$ and the true $z_+$ tends to zero.
On one fixed small coefficient box, \eqref{eq:counting-connector}
therefore has a transfer analytic on a slightly larger box and with
state room.  Its scalar matching equation is this transfer from $0$
minus $(Y_+-Y_-)/r$.

A word records the passage cells, real charts, signs, traversal
orientations, and endpoint-rule alternatives of
Theorem~\ref{thm:terminal-reduction}, together with these connector
choices.  Each choice is finite and the word length is bounded, so
$\mathcal W_d$ is finite.  Internal subdivisions used only in an
absolute zero test are not part of $\mathcal W_d$.

\paragraph{Parameters, endpoints, and argument graphs.}
Use the polynomial coefficients and all parameter-only preparation
data for the selected word as external coordinates $p_w$.  Preparation
was performed with the field coefficients as parameters, before
choosing orbit cuts.  Thus, for a fixed field and fixed preparation
choices, these values are fixed for every cycle assigned that word.
They can vary independently on the extended analytic parameter
domain; no global analytic relation between different parameter-only
preparation formulas is required there.  Any quantity depending on
the endpoints, including a normalization by a proposed endpoint or a
connector coefficient in \eqref{eq:counting-connector}, is instead an
unknown determined by its argument graph.

For a word with $k$ links introduce physical endpoints
$e_j=(x_j,y_j)$, $j\in\ZZ/k\ZZ$, shared by successive links.
There is no section equation fixing the tangential position of a cut.
The other variables are terminal coordinates, clocks, amplitudes,
bounded ordinary arguments, and residuals.  We use the following
encoding conventions.
\begin{enumerate}
\item Earlier time-dependent coordinate formulas needed to recover
physical endpoints are included in each clock preparation.  Their
values are positive rational clock powers, products with
parameter-only coefficients, and bounded analytic units at prepared
arguments.  If $q>0$ and $\nu=A/B\in\mathbb Q$, $B>0$, the equation
$u^B=q^A$, on $u>0$, defines the selected positive power regularly;
negative $A$ is handled by a nonzero denominator.  Translations,
products, and ratios have their usual regular graph equations.
Bounded analytic state changes, simple-root linearizers, diagonal
primitives, and their regular inversions use smaller boxes with
extensions to larger neighborhoods.  Complex conjugate preparation
data are paired and expressed through real and imaginary parts, so
all equations here are real analytic in real coordinates.
\item A logarithmic clock length is tied by
\begin{equation}\label{eq:counting-logties}
 t_-=t_+e^{-L},\qquad L>0,\qquad 0<t_-<t_+.
\end{equation}
Its derivative in $L$ has absolute value $t_-$ and is nonzero.
For an additive clock we additionally use
\[
 q=pe^{-L_w},\quad L_w>0,\quad
 Q=p+\delta qL_w,\quad S=Q-q.
\]
Products tie scaled actions such as $W$, slopes, and bounded
prefactors.  A logarithmic state residual is tied using the positive
physical transformed state $R=e^v$:
\begin{equation}\label{eq:counting-residualtie}
 R_+=R_-\exp(\kappa S+z_+),
\end{equation}
or the corresponding equation with $S=W$ or zero subtracted slope.
Place the decaying exponential on the appropriate side when useful.
The equation has nonzero derivative in $z_+$, and hence is a regular
graph.  Absolute unbounded state logarithms do not become analytic
arguments.  Every unbounded exponential in these ties is a fixed
affine combination of clock coordinates; a variable scaled action
has its own coordinate and product tie.
\item Right amplitudes contain their prescribed bounded residual
correction.  For example in an additive annulus, if
$z_+=v_+-v_--\kappa S$, then
\[
 D_i=\chi_iQ^{p_i}e^{d_i v_+-d_i z_+}
     =\chi_iQ^{p_i}e^{d_i v_-+d_i\kappa S}.
\]
Thus the proposed final state cancels from the actual-flow amplitude.
The analogous mixed identity is
\[
 \chi_i e^{a_i\theta_++\beta_i v_+-\beta_i z_+}
 =\chi_i e^{a_i\theta_++\beta_i v_-+\beta_i\kappa W}.
\]
These are identities on the tied endpoint graph, including mismatched
endpoints.  Signed or zero amplitudes are handled by products, without
extracting roots from their signs.  An invariant diagonal expression
is evaluated at one end only.  Its weight-zero identity makes it
constant under the relevant nearby endpoint variation at fixed field
parameters.  No redundant second equality is added for that identity.
\item Each link has exactly one scalar transfer equation in addition
to its argument ties.  A mixed kernel uses independent normalized
clock coordinates $(L,W)$, with exponents $\alpha L+\beta W$ for
fixed $\alpha,\beta$.  Its other bounded slope or coefficient copies
are separate ordinary arguments with their own ties.  The additive
kernels have independent allowed $(Q,q,\delta)$ and layer data.
Duplicate arguments may be given duplicate coordinates, with regular
equality ties.  All kernels are the actual real analytic ODE
transfers of the preceding sections.  The residual extensions with
independent arguments are those of
Lemma~\ref{lem:terminal-ambient-extension}, preserving the required
layer-vanishing identities.
\end{enumerate}
Each auxiliary coordinate has one regular determination in this
construction.  Ordered in the construction order, the graph
Jacobian is block triangular with invertible diagonal.  In
particular its solutions are locally unique graphs over physical
endpoints and fixed parameters.

Figure~\ref{fig:matching-arguments} separates the independent kernel,
the tied endpoint graph, and its matching locus.  Only the last step
imposes the scalar transfer equation.
\begin{figure}[htbp]
\centering
\begin{tikzpicture}[
  box/.style={draw, rounded corners=2pt, align=center,
    text width=3.9cm, minimum height=3.2cm, inner sep=5pt,
    font=\small},
  flow/.style={-{Stealth[length=2mm]}, thick},
  note/.style={font=\footnotesize, align=center}
]
\node[box] (ambient) at (0,0) {
  \textbf{Independent arguments}\\[4pt]
  $a\in U$\\[3pt]
  $\Phi(a),\quad D_a\Phi$
};
\node[box] (graph) at (5.9,0) {
  \textbf{Proposed endpoints}\\[-1pt]
  \textbf{with argument ties}\\[4pt]
  $e=(e_-,e_+),\quad a=A(e)$\\[3pt]
  $\mathcal M(e)=\Phi(A(e))-z_+(e)$\\[3pt]
  {\footnotesize Nearby endpoints may mismatch}
};
\node[box] (matched) at (11.8,0) {
  \textbf{Matching endpoints}\\[4pt]
  $a=A(e)$\\[3pt]
  $\mathcal M(e)=0$
};
\draw[flow] (ambient.east) -- node[note, above=3pt]
  {impose\\$a=A(e)$} (graph.west);
\draw[flow] (graph.east) -- node[note, above=3pt]
  {impose\\$\mathcal M=0$} (matched.west);
\node[note] at (5.9,-2.05) {
  $D_e\mathcal M=(D_a\Phi)(A(e))\,D_eA-D_ez_+$
};
\end{tikzpicture}
\caption{One link near an actual passage, with physical parameters fixed.
The arrows impose successive equations.  Regular argument ties express
all independent kernel arguments, including coefficient copies, in terms
of nearby proposed endpoints before the scalar transfer equation is
imposed.  Here $z_+(e)$ is the prescribed output coordinate in the chosen
traversal.  On this local tied graph, $\mathcal M$ is the physical mismatch
in regular transverse coordinates, including at mismatched pairs.
Ambient partial derivatives are taken in $a$; endpoint derivatives use
the displayed chain rule.  This is a local statement, not a claim that
every point of the extended matching domain is a physical passage.}
\label{fig:matching-arguments}
\end{figure}


\paragraph{The specified open domains.}
We list all domain restrictions as strict margins: positive clock
and root conditions; ordering; all nonzero divisors; coordinate-box
and smallness restrictions for bounded analytic arguments; coefficient
signs and ranges; and the state and rate margins ensuring that the
real transfer exists in its larger state disk.  The finite analytic
charts are chosen with a larger extension about their closed smaller
argument boxes.  The target smallness thresholds from the geometric
reduction are chosen strictly inside these boxes.

For additive transfers this includes $Q>q>q_*$ for the fixed large
threshold $q_*$, a sufficiently small $|\delta|$, and small layer
amplitudes and state inputs.  The action is increasing.  The real
estimates of Theorem~\ref{thm:additive-packets} give state room:
model (A) has damping and a small integral of the nonconstant-power
forcing; in model (B) the large $q_*$ makes the adverse fixed power
factors harmless in the integrals of the fast layers.  For mixed
transfers, fixed rate-sign margins, small amplitudes, and smaller
state boxes give the corresponding room from
Theorem~\ref{thm:mixed-packets}; stable orientation is used in the
damped model.  The regular and bounded-action alternatives use the
integrated bounds of Theorems~\ref{thm:regular-packets} and
\ref{thm:terminal-reduction}.  For example their small integrated
coefficients are explicitly tied arguments of the form $BL$,
$B\tau_e^\lambda$ on a long-length alternative, or
$LB\tau_e^\lambda$ on a short-length alternative.  Alternative
length ranges overlap about their cutoffs.  These finite conditions
are coordinate, polynomial/product, or bounded analytic inequalities,
with the same affine clock exponentials already allowed in the
argument graphs.

Thus a bounded sequence whose listed margins stay bounded away from
zero has a limit in the same boxes, sign and ordering domains, and
real existence regime.  Every formula is analytic at that limit and
the limit is admissible.  No additional condition that a long
mismatched arc remain in a physical subdivision cell is imposed.
The analytic systems are the normalized kernels on the stated
domains; their identification with physical motion is needed locally
at the target representations.  This distinction supplies exactly the
domain completeness required in Definition~\ref{def:counting-domain}.
It also remains valid after adding a determinant margin: a bounded
sequence with that margin bounded below cannot leave its minor
chart.

\paragraph{Local physical validity and rank.}
At a target representation and its fixed actual parameter value,
Lemma~\ref{lem:terminal-tied-endpoints} identifies each tied transfer
equation with a true scalar-flow endpoint mismatch, up to a regular
state coordinate and a nonzero factor.  The connector has the same
property by \eqref{eq:counting-connector}.  The subarcs are interior
and their clocks have nonzero velocity, so this identity holds for
all nearby proposed endpoints, not only for matching pairs.

One can see directly why a normalization chosen from the proposed
endpoint does not spoil the differential.  For a held normalization
$a$, let $H(e_-,e_+,a)$ be the transformed mismatch.  In local
coordinates consisting of a physical scalar mismatch $\Delta_0$
and the remaining variables, regularity gives
$H=U(e_-,e_+,a)\Delta_0$ with $U\ne0$.  This follows, for example,
by integrating $\partial_{\Delta_0}H$ from $0$ to $\Delta_0$.
After substituting $a=a(e_-,e_+)$, its differential on
$\Delta_0=0$ is still $U\,\dd\Delta_0$.
If the terminal clock is also transformed, its nonzero clock
velocity determines the hitting-time variation; the regular
coordinate pair leaves a nonzero transverse mismatch.  The exact
residual-amplitude identities displayed above and the weight-zero
invariance of the diagonal terms ensure the premise for all the
terminal transformations used here.

Choose local flow coordinates $(s_j,r_j)$ at the $k$ endpoints,
with $s_j$ tangential and $r_j$ transverse.  Up to nonzero row
factors, the link differentials at a matching cycle are
\begin{equation}\label{eq:counting-cyclic-rank}
 \dd r_{j+1}-h_j\,\dd r_j,
 \qquad j\in\ZZ/k\ZZ,
\end{equation}
where $h_j\ne0$ is the derivative of the corresponding section
holonomy.  Reversing the orientation used to write one link merely
changes this equivalent equation by a nonzero factor.  The kernel
of the cyclic transverse matrix satisfies
$\dd r_1=(\prod_j h_j)\dd r_1$.  For a hyperbolic orbit its return
multiplier is $\prod_j h_j\ne1$, so that kernel is zero and the
$k$ link rows have rank $k$.  The tangential variables remain free,
as intended.  Adding the auxiliary graph equations adds an
invertible block in their own variables.  Equivalently, eliminating
that block by row operations leaves exactly
\eqref{eq:counting-cyclic-rank}.  Hence the complete equation
Jacobian has full row rank, and at least one of its finitely many
maximal minors is nonzero.

For a word with $k$ links and $n_{\rm aux}$ auxiliary unknowns, the
fiber has $2k+n_{\rm aux}$ unknowns and $k+n_{\rm aux}$ equations.
Near a hyperbolic target its dimension is therefore $k$, accounted
for by the tangential positions of the cuts.  Hyperbolicity gives
full row rank of this matching system.  The following argument shows
that a connected component cannot contain targets from two different
cycles.

\paragraph{Injection into analytic components.}
Fix a word, its actual parameter tuple, and one nonzero-minor chart.
Let $M$ be a connected component of its analytic fiber and let
$E:M\longrightarrow\RR^2$ map a representation to its first
physical endpoint.  Suppose $M$ contains a target on a hyperbolic
cycle $C$.  In a neighborhood of that representation, exact local
matching and uniqueness of the nearby fixed point of the return
map imply $E\in C$.  The free tangential cuts only move points
along $C$.

We claim that $E(M)\subset C$.  The periodic trajectory $C$ is a
compact embedded nonsingular real analytic curve: it is analytic
by analytic ODE dependence, regular because $V$ has no zero on a
nonconstant periodic trajectory, and embedded by uniqueness and
use of a minimal period.  Put
$A=\operatorname{int}_M E^{-1}(C)$.  This is nonempty and open.
If $q\in\overline A$, closedness of $C$ gives $E(q)\in C$.
Near $E(q)$ a regular analytic defining function $h$ cuts out
exactly $C$.  Choose a connected analytic neighborhood $B$ of $q$
with $E(B)$ in that neighborhood.  The analytic function $h\circ E$
vanishes on the nonempty open set $A\cap B$, so the analytic
identity theorem makes it zero on $B$.  Thus $q\in A$; $A$ is
closed, and connectedness gives $A=M$.

Different periodic trajectories are disjoint by uniqueness of the
ODE.  They therefore cannot have target representations on the
same component $M$.  This argument permits remote points of $M$
that are not physical passage representations.  It uses physical
validity on the initial neighborhood and analytic continuation on
the entire component, which is precisely why no additional global
physical-cell restriction was imposed.

For each word we have now constructed a fixed open analytic system.
Its domains are positive-margin complete, every retained cycle has a
representation, and different hyperbolic cycles occupy different
components on each maximal-minor chart.  To apply the projection theorem,
it remains to prove absolute isolated-zero finiteness for every finite
system generated from these data.

\subsection{Absolute zeros of the matching systems}

Fix the matching systems just constructed, including their independent
kernel arguments, graph ties, and strict domain restrictions.  A system
in the counting closure may also contain derivatives, inverse margins,
and multiplier variables.  When its arguments approach a boundary, we
will rewrite its kernels by exact identities and retain every added
graph equation.  The next two lemmas state what these rewritings preserve;
the recovery lemma records the consequence used when a free flag input
is varied in the final contradiction.

\begin{lemma}[Graph and finite-jet lifting]\label{lem:counting-jet-lift}
Suppose an analytic transfer $\Phi(u)$, on an open set of its original
independent arguments $u$, has an exact analytic decomposition into
finitely many transfers, with auxiliary arguments and intermediate
states determined by regular, locally unique graph equations.  A
finite system involving $\Phi$ and finitely many of its independent
derivatives can be replaced, locally at any admissible tuple, by a
system involving those transfers and finitely many of their
derivatives.  The replacement, with the graph equations retained,
is locally analytically isomorphic to the original solution set.
In particular it preserves isolated solutions.  This conclusion also
holds along a sequence with shrinking graph neighborhoods, provided
one fixed finite decomposition is used along the sequence.
\end{lemma}

\begin{proof}
If $G(u,v)=0$ has $G_v$ invertible, its selected solution is a
locally unique analytic graph $v=v(u)$.  Projection of that graph to
$u$ is an analytic isomorphism.  Substituting $v(u)$ in any other
equations proves the first assertion about their solution sets.
Denominator clearing preserves this assertion on the retained
nonzero domain.

For the derivative assertion, differentiate the identity expressing
$\Phi$ in terms of the other transfers \emph{after} their arguments
have been put on their graphs.  Implicit derivatives of $v(u)$ are
obtained successively from $G_v^{-1}$ and finite derivatives of $G$.
The finite chain rule expresses each requested derivative of $\Phi$
using only finitely many such quantities.  They may either be
substituted, or be given their own graph coordinates.  In the latter
description, the equations for derivatives of order $j$ are linear
with coefficient $G_v$ in the order-$j$ unknowns and have previously
determined lower-order terms.  Their complete finite-jet Jacobian is
block triangular with invertible diagonal.

In particular, for a subdivision into $K$ consecutive subflows,
write
\[
 z_j=\Phi_j(a_j(u),z_{j-1}),\qquad 1\le j<K,
 \qquad
 \Phi(u)=\Phi_K(a_K(u),z_{K-1}).
\]
The intermediate-state equations have a triangular Jacobian with
unit diagonal.  If jets are adjoined, each derivative of $z_j$ is
determined from the earlier $z_i$ and their jets; these graph equations
again have unit diagonal.  Thus the augmented system has precisely
the original germ of solutions and no additional local degree of
freedom.  This is a pointwise assertion at each member of a sequence;
no uniform radius of the isomorphism is needed to preserve isolation.
\end{proof}

Two consequences will be used below.  First, partial derivatives of
two independently enlarged kernels cannot be identified just because
their values agree on a graph.  Only the differentiated composition
identity in Lemma~\ref{lem:counting-jet-lift} is used.  Second, a
finite $K$ chosen after taking a subsequence is legitimate in an
absolute zero contradiction: it produces one fixed finite system for
that sequence.  It does not assert a degree-uniform value of $K$ and
does not add links to the geometric word count.

\begin{lemma}[Cleared finite differential closure]
\label{lem:counting-cleared-closure}
Consider one fixed finite collection of primitive functions, their
argument graphs, and a finite differential and algebraic construction
from them.  Assume that, after one fixed exact graph enlargement, the
primitive functions and defining graph equations extend analytically
near a finite limiting tuple.  The selected graph Jacobians and every
declared divisor need be nonzero only at the points of the sequence,
not at its limit.  Then the constructed equations can be replaced by
analytic equations near that limit, retaining all graph equations and
all original open restrictions, with exactly the same solution germ
at each selected point.  Additional multiplier coordinates, polynomial
equations, and parameter copies or specializations preserve this
assertion.
\end{lemma}

\begin{proof}
First differentiate every exact rewriting as an identity on its old
open argument domain.  For a graph $G(u,v(u))=0$ this gives
\[
 v_{u_i}=-G_v^{-1}G_{u_i}.
\]
Cramer's rule and induction on derivative order express every requested
implicit derivative as a quotient of polynomials in finite jets of
$G$, with denominator a power of $\det G_v$.  The same argument
applies successively to a finite stack of graphs and to the finite
composition identities in Lemma~\ref{lem:counting-jet-lift}.

Every further algebraic or differential operation preserves a rational
expression in a finite collection of analytic jets.  For example,
\[
 \partial_i(A/B)
   =\frac{B\partial_iA-A\partial_iB}{B^2},
 \qquad (A/B)^{-1}=B/A
\]
on the retained nonzero domains.  An inverse of a new expression
therefore adds its nonzero numerator to the possible denominator
factors; a determinant is a polynomial in its entries.  Thus each
final equation is $A/B=0$, where $A,B$ are analytic near the limiting
tuple and $B$ is nonzero at every selected point.  Replacing it by
$A=0$ leaves its solution germ unchanged there.  This replacement
does not require a bound on $B^{-1}$, nor does it append its value
as a new coordinate.  In particular a margin or a graph determinant
tending to zero causes no difficulty at this step.  Finite jet
coordinates, if retained instead of substituted, have the
differentiated analytic graph equations of
Lemma~\ref{lem:counting-jet-lift}; their Jacobians need be invertible
only at the selected points.

New multipliers are coordinate factors in these expressions.  A
parameter copy has a polynomial equation $p'-p=0$, and a fixed
specialization has an equation $p-p_0=0$.  Unrestricted added
coordinates likewise enter the closure through polynomial and
rational operations.  These operations introduce no new primitive
function or evaluation outside its original argument domain.  Their
finite derivatives are of the same form, proving the last assertion.
\end{proof}

\begin{lemma}[Recovery after graph enlargement]
\label{lem:counting-graph-recovery}
Let $u_i$ be a sequence of isolated solutions of a fixed system,
where $u$ denotes its full tuple, including parameters and multipliers.
Suppose an exact finite enlargement appends $v_i$ by equations
$G(u,v)=0$ whose selected graph Jacobian $G_v(u_i,v_i)$ is invertible
at every point.  Retain all these equations.  An old flag input
recoverable from $(u,v)$ cannot vary along the local lifted solution
set while $u$ stays fixed, provided the variation is taken within the
selected graph neighborhood.  These neighborhoods may shrink
arbitrarily along the sequence.
\end{lemma}

\begin{proof}
At each $i$ the implicit function theorem supplies an open neighborhood
$N_i$ in which $v=g_i(u)$ is the unique selected solution of the
complete graph equations.  Let $B_i$ be any proposed isolation
neighborhood of $u_i$.  Intersect $N_i\cap\pi^{-1}(B_i)$, where
$\pi(u,v)=u$, with the local real flag chart and all strict domain
neighborhoods at that point.  This is an open neighborhood of the
lifted point.  By continuity of the chart map, all sufficiently small
variations of a free chart input remain in it.  The size is chosen
separately for every $i$; it can also be made small enough to preserve
the prescribed comparisons and ordinary limits for the retest in
Theorem~\ref{thm:absolute-zero}.

Write $S_*=R(u,v)$ for the retained old free input, recovered by the
locally invertible affine and lifted logarithmic flag changes.
For any such varied solution, $u'=u_i$ would imply
$v'=g_i(u_i)=v_i$ and hence $S_*'=S_*$, a contradiction if this
input was changed.  Thus a nonzero sufficiently small permitted
variation gives a different original tuple inside $B_i$.  Vanishing
graph determinants at the limiting tuple affect only the sizes of
the individual neighborhoods, not this argument.
\end{proof}

\begin{proof}[Completion of the proof of Theorem~\ref{thm:matching-system}]
Fix an arbitrary finite system in the closure required by
Theorem~\ref{thm:projection-count}, including its equations and every
original graph and strict domain restriction.  Suppose it has a
sequence of distinct isolated total solutions.  Multipliers and
coordinates formerly called parameters are now ordinary unknowns.
Clear all declared nonzero denominators, retaining their restrictions.
Only finitely many derivatives of finitely many transfer kernels,
argument graphs, and margins occur, together with algebraic factors
in the coordinates.
This assertion is hereditary under the projection construction:
gradients of inverse margins, products with new multiplier coordinates,
and determinants of their derivatives have exactly this form by
Lemma~\ref{lem:counting-cleared-closure}.  Parameter copies and
specializations add polynomial equalities.  Thus an unbounded new
coordinate is included in the large-coordinate analysis below, while
a bounded new coordinate is just an additional ordinary argument of
these polynomial expressions.  No new transfer family is introduced.

At every transfer choose a subsequence regime of
Theorem~\ref{thm:terminal-reduction}; there are finitely many transfers
in this fixed system.  Append any bounded ratios or auxiliary
coordinates required by that regime by their unique regular graph
ties.  For example, the bounded-ratio additive corner is normalized
by
\begin{equation}\label{eq:counting-ratio}
 S=Q-q,\qquad M=S/q,\qquad
 y+\delta\log y=1+Ms.
\end{equation}
If $M$ tends to a finite value, a fixed sufficiently large finite
subdivision into $K$ fractions gives regular analytic normalized
pieces.  The value of $K$ is chosen once for this limiting corner.
Its internal states and amplitudes are regular graphs.  Internal
amplitudes are products of bounded positive ratios and exponentials
of fixed affine clock fractions; for instance a left additive layer
at fraction $j/K$ has amplitude
\[
 e_j=e\,(w_j/q)^\beta e^{-a jS/K}.
\]
The real layer estimates keep these amplitudes in the permitted
boxes.  The mixed unbalanced auxiliary ratio $C_{\rm aux}=P/R$ is
also appended where required by the mixed packet construction; here
$P,R$ denote that construction's positive scales, not the components
of the polynomial vector field.

All derivatives of the original kernel are obtained by differentiating
the exact identities on its original open argument domain.  For
example if the decomposition in \eqref{eq:counting-ratio} reads
$\Phi(Q,q,\delta,a)=\Psi(S,M,\delta,a)$, then
\[
 \Phi_Q=\Psi_S+q^{-1}\Psi_M,
 \qquad
 \Phi_q=-\Psi_S-(Q/q^2)\Psi_M.
\]
Higher derivatives use finitely many rational factors and kernel
derivatives, with the same nonzero denominators.  The other bounded
clock and subdivision identities are handled by the same chain rule,
as proved in Lemma~\ref{lem:terminal-ambient-extension}.
Lemma~\ref{lem:counting-jet-lift} therefore lifts the isolated
solutions, including any equations containing derivatives, to
isolated solutions of one fixed finite enlarged system.  Using a
different $K$ at every term of the sequence would not give this
conclusion; the fixed finite choice just described does.

\paragraph{Bounded enlarged tuples.}
Suppose first that every coordinate of the enlarged tuple remains
bounded, and pass to a finite limit.  Some strict margins may vanish.
Positive-margin domain completeness, used for the proper barrier,
does not prove analytic extension at such a limit.
Lemma~\ref{lem:terminal-boundary-corners} instead gives the required
extension of each \emph{old independent-argument} normalized kernel:
closed smaller argument and amplitude boxes lie within larger analytic
boxes, and the integrated-field bounds leave a positive state margin
on their closures.  The bounds use independent coefficient copies,
before any physical slope or normalization ties are imposed.  At a
bounded clock corner these normalized fields therefore define analytic
transfers near the limiting data, including when a clock is zero.

For example, the normalized mixed equations have the forms
\[
 v_s=-Wv+L g(X,Y,v),\qquad z_s=L G(X,Y,z),
\]
with layer exponents $\alpha L+\beta W$ and other coefficient copies
retained independently.  They contain no singular $W/L$ in these
positions and are analytic also at $(L,W)=(0,0)$.  A slope needed
elsewhere is its own coordinate, tied by $L c-W=0$.  If it stays
bounded, this polynomial graph is analytic at the corner.  If a
necessary new ratio diverges, the enlarged tuple has a large coordinate
and belongs to the large-coordinate case below; boundedness of the
\emph{old} tuple alone is not the criterion for the present case.

The remaining possible degenerations of the endpoint graphs are
equally explicit.  A ratio has an equation $bv-a=0$.  A positive
rational root has an equation $u^B-q^A=0$, made polynomial by
clearing negative powers when necessary.  Clock and residual ties have
the forms
\[
 t_- -t_+e^{-L}=0,
 \qquad R_+-R_-e^{\kappa S+z_+}=0.
\]
They extend analytically when all their coordinates remain finite,
even if their derivatives in $L$ or $z_+$ tend to zero with the
positive endpoints.  If $t_+/t_-\to\infty$, then
$L\to\infty$, which is a large-coordinate case.  Connector
coefficients have polynomial graph equations after multiplication by
$V_X(X_-,Y_-)$ and $r$, and $r^2=\Delta$ is polynomial.  Their
bounded normalized coefficient arguments remain in the larger
analytic connector box when either divisor tends to zero.
Regular endpoint inversions may likewise retain their analytic
defining equations.  Their graph Jacobians need only be nonzero at
the selected points.  Finite implicit derivatives then have exactly
the denominator description of
Lemma~\ref{lem:counting-cleared-closure}.

Consequently all primitive transfers and defining graph equations
extend analytically at the finite limiting tuple of the enlarged
system.  That lemma clears
the finite derivative and multiplier equations without asserting
invertibility of a limiting graph Jacobian.  Only equations are
cleared; no inequality is replaced by an unverified numerator sign.
The old strict inequalities and all later determinant restrictions
remain as open restrictions at each selected point.  Their limiting
values may vanish, and their differentiated forms need not themselves
extend across that boundary.  This is sufficient here, in contrast
to the positive-margin completeness used in the projection theorem.
Local finiteness of components
of a real analytic set excludes an accumulating sequence of its
isolated zeros \cite[Corollary~2.7]{BierstoneMilman1988}.
Each point of the original strict domain is interior to that domain,
so it cannot become an isolated zero merely by intersecting a
nonisolated germ with those open restrictions.  This excludes the sequence of isolated solutions when its enlarged
tuple has a finite limit.

\paragraph{Divergent enlarged tuples.}
Suppose instead that at least one enlarged coordinate diverges.
Split all coordinates of the enlarged sequence into bounded
convergent coordinates and signed large coordinates.  Apply the
affine flag construction of Section~\ref{sec:flags} to the latter,
keeping bounded remainders as ordinary coordinates.  Convert the
initial basis scales, including scales used in polynomial factors,
to dependent logarithmic nodes, and insert the further logarithms
required by the primitive constructions.  Saturate the resulting
finite flag.  These real coordinate changes are locally invertible.
The kernel rates are affine in the primary clock coordinates and
bounded remainders.  Polynomial powers of the initial scales are
now shift monomials times bounded analytic factors.  Under an
inserted logarithm, old independent derivatives become finite sums
of packet independent derivatives multiplied by constants, signs,
and reciprocal powers of the logged dependent scale.  These are
precisely operations of Theorem~\ref{thm:calculus}.
If only a multiplier or an algebraic auxiliary coordinate diverges
while a kernel's clocks remain bounded, that kernel keeps its
ordinary analytic normalized formula.  It is a packet constant in
the extra free inputs before pullback; the new polynomial coordinate
factors are treated by the preceding logarithmic conversion.  In
particular a large auxiliary ratio at bounded primary clocks does
not invoke an unbalanced mixed theorem that requires both clocks
large.

Theorem~\ref{thm:terminal-reduction}, with the differentiated packet
theorems, supplies the packet data for this finite derivative
request on independent flag inputs.  Graph ties remain equations;
they are not assumed to hold when taking independent tests of a
primitive kernel.  In particular a bounded coefficient copy and a
clock ratio remain separate arguments until their graph identity is
used by the chain rule.  The layer-vanishing residual extensions of
Lemma~\ref{lem:terminal-ambient-extension} are essential here.
For each nonzero equation, iterated left-finiteness permits a leading
shift normalization making both sign trees nonnegative in the series
sense.  A formally zero equation is handled by separation.  All
normalizing monomials are nonzero on the real domain and use the
current flag backgrounds.  The finitely many real normalized
equations can then be combined by a sum of squares, exactly as in
Theorem~\ref{thm:absolute-zero}; the resulting positive operations
are covered by Theorem~\ref{thm:calculus}.

\paragraph{Retesting and return to the original solutions.}
We verify the retestability of this assembled library in the precise
sense of Definition~\ref{def:elimination-retestable}.  For the
primitive transfers this is the coefficient-germ assertion of
Theorems~\ref{thm:additive-packets}, \ref{thm:regular-packets}, and
\ref{thm:mixed-packets}, carried through the exact terminal identities
by Theorem~\ref{thm:terminal-reduction}.  The affine and logarithmic
flag identities, fixed analytic fields, sector signs, and ordinary germs
are unchanged by sufficiently small absolute perturbations preserving
the selected comparisons.  Terminal coefficients are the fixed
formal recursions and, where an infinite anchor is used, its fixed
normalized jets.  Cutoff and finite-anchor choices have already
dispersed before these terminal leaves are taken.  Subsequent
algebra and finite derivatives preserve those same defining germ
identities.  The implicit inversion of
Theorem~\ref{thm:calculus} is used only when the limiting
exponent-zero Jacobian is invertible.  A root, logarithmic, or ratio
graph that is regular at every selected point but degenerates at the
limit is instead retained as an equation in its own coordinates;
it is not inverted by the packet implicit theorem.  Its finitely
many old-coordinate derivatives are expressed by Cramer's rule and
their nonzero denominators are cleared, leaving finite algebraic
expressions in independent jets of the defining equations.  At
large scales these equations use the already admitted affine shifts
and bounded analytic fields.  The ordinary elimination theorem
then treats them together with all the other equations, including
any multiple limiting zero.  Thus only actual unit-Jacobian
ordinary chart operations invoke the implicit part of
Theorem~\ref{thm:calculus}; no packet inverse at a degenerate
limiting graph is assumed.  Further flag refinements leave the old
free inputs free.  At least one such input is recoverable from the
tuple before the flag changes whenever a large scale remains.  The
sequence-wise graph enlargement retained all its defining equations,
so Lemma~\ref{lem:counting-graph-recovery} makes it recoverable for
the original isolation test as well.  More explicitly, at each
anchor choose the free variation inside that anchor's graph-uniqueness
neighborhood, the real flag chart, and the pullback of the proposed
original isolation neighborhood.  Such variations may be chosen
arbitrarily small, with no common radius required along the sequence.
All the hypotheses of
Theorem~\ref{thm:absolute-zero} therefore hold.  It excludes the
isolated sequence in that case; the small free variations used there
can be made inside every strict open condition at the selected
point.

The two alternatives prove the absolute zero property for the arbitrary
finite system under consideration.  All additional determinant margins
and multiplier systems use the same finite operations, so the conclusion
holds for them as well.

Together with the construction and geometric arguments above, this
proves all three assertions of Theorem~\ref{thm:matching-system}, including
the assertion about maximal-minor charts.
\end{proof}

\subsection{Rotation and the bound in the whole plane}

\begin{lemma}[Hyperbolic realization of a finite cycle count]
\label{lem:counting-rotation}
For any finite collection of limit cycles of a planar polynomial
field $V=(P,Q)$, a sufficiently small perturbation
$V_\mu=V+\mu(-Q,P)$, of one common nonzero sign, has at least as
many distinct hyperbolic periodic orbits in the union of disjoint
neighborhoods of that collection.  The degree bound is preserved.
\end{lemma}

\begin{proof}
Let $z(t)$ parametrize one selected cycle with period $T$, and start
the perturbed solution at the same section point.  Set
$u(t)=\partial_\mu z_\mu(t)|_{\mu=0}$ and
$J(a,b)=(-b,a)$.  Then
\[
 u'=DV(z(t))u+JV(z(t)),\qquad u(0)=0.
\]
For $w(t)=\det(V(z(t)),u(t))$, the identity
$\det(Aa,b)+\det(a,Ab)=\operatorname{tr}(A)\det(a,b)$ gives
\begin{equation}\label{eq:counting-rotation}
 w'=(\operatorname{div}V)(z(t))w+|V(z(t))|^2,
 \qquad
 w(T)=\int_0^T
 e^{\int_s^T(\operatorname{div}V)(z(\tau))\,\dd\tau}
 |V(z(s))|^2\,\dd s>0.
\end{equation}
The return-time correction adds a multiple of $V$ to the terminal
variation and does not change this determinant.  Thus the local
analytic Poincare displacement $D(r,\mu)$ has
$D_\mu(0,0)\ne0$ in a transverse section coordinate $r$ centered
at the cycle.

The analytic implicit function theorem writes its zero set as
$\mu=\psi(r)$ near $(0,0)$, with $\psi(0)=0$.  Isolation of the
cycle makes $D(r,0)$ a nonzero analytic germ, so
\[
 \psi(r)=a r^m+O(r^{m+1}),\qquad a\ne0,\quad m\ge1.
\]
For odd $m$, both sufficiently small signs of $\mu$ give one
nearby nonzero root; for even $m$, the sign of $a$ gives two.
All these roots are simple because
$\psi'(r)=ma r^{m-1}+O(r^m)\ne0$ for sufficiently small
$r\ne0$.  Differentiating $D(r,\psi(r))=0$ then gives
$D_r=-D_\mu\psi'\ne0$ at these roots.  A planar return map has
positive derivative (as follows from its scalar variational
equation), so a simple fixed point has multiplier different from
$1$ and is hyperbolic.

Choose disjoint tubular neighborhoods of the finitely many original
cycles and one common sufficiently small magnitude of $\mu$.
If $O$ cycles have odd order and $E_+$ and $E_-$ have even order
favoring the two signs, the resulting numbers are at least
$O+2E_+$ and $O+2E_-$.  Their maximum is at least
$O+E_++E_-$, the original selection size.  The count is taken over
the union: an even-order cycle favoring the other sign need not
persist for the chosen rotation.  The new cycles are
distinct because their neighborhoods are disjoint (and the two
roots in one sufficiently small return section give different
periodic orbits).  Finally $V_\mu$ is obtained by constant linear
combinations of $P,Q$, and still has degree at most $d$.
\end{proof}

\begin{proof}[Proof of Theorem~\ref{thm:uniform}]
Fix $d$.  For every word $w$ of Theorem~\ref{thm:matching-system}
and every one of its maximal-minor charts $I$, apply
Theorem~\ref{thm:projection-count}.  Let $N_{w,I}<\infty$ be the
resulting component bound.  It is independent of the numerical
value of every polynomial coefficient and preparation parameter.
The component injection gives, in the fixed square, the uniform
bound
\begin{equation}\label{eq:counting-final-bound}
 B_*(d)=T(d)+\sum_{w\in\mathcal W_d}\ \sum_I N_{w,I}
\end{equation}
for hyperbolic periodic orbits.  An orbit represented in more than
one chart is assigned to any one of them; overcounting in this sum
is harmless.  Finiteness of the sums is geometric template
finiteness, and does not involve the internal $K$ used in a
sequence-dependent absolute zero test.

Now select any finite collection of limit cycles of any field of
degree at most $d$ in the whole plane.  A spatial affine scaling
$z=a+Ru$, $R>0$, puts all of them strictly inside the fixed square.
The transformed field is
$\dot u=R^{-1}V(a+Ru)$, again polynomial of degree at most $d$;
isolation and geometric distinctness are preserved.  Apply
Lemma~\ref{lem:counting-rotation} there.  The perturbation can be
chosen small enough that its retained hyperbolic cycles remain in
the square.  Their number is at least the size of the selected
collection and is at most $B_*(d)$ by
\eqref{eq:counting-final-bound}.  Hence every finite collection has
size at most $B_*(d)$.  If the original field had more limit cycles,
one could select $B_*(d)+1$ of them, a contradiction.  Taking any
integer upper bound for $B_*(d)$ defines the required finite
$B(d)$, without an assumption of prior individual finiteness.
\end{proof}
